\chapter{Abstract Reduction and String Rewriting}
\label{ch:abstract-string-rewriting}

Deleting redundant separators looks like a harmless textual cleanup until two rules disagree about which separator survives.
An optimizer faces the same structural problem when one rewrite exposes another and a different choice prevents it.
Before choosing a rule engine, we need vocabulary for possible steps, final answers, and the role of strategy.
This vocabulary applies to strings, syntax trees, shared graphs, and compiler states without pretending that they are interchangeable representations.

\section{Separate the relation from the strategy}
\label{sec:ch81-reduction}
An \term{abstract reduction system} consists of a carrier set $A$ and a binary one-step relation $\to$ on $A$.
We write $a\to^*b$ for zero or more steps, $a\to^+b$ for one or more, and $a\leftrightarrow^*b$ for conversion allowing steps in either direction.
A \term{normal form} has no outgoing step.
A reduction strategy selects among the available steps; the relation describes every legal choice, while a particular implementation may explore only one path.

For strings over $\Sigma$, a rule $l\to r$ permits $ulv\to urv$ for arbitrary surrounding strings $u,v$.
This closure under context is what turns a finite rule collection into a relation on infinitely many strings.
An occurrence is described by a start offset and matched extent, measured in the chosen alphabet.
If the alphabet is tokens, a rule cannot start inside a token; if it is bytes, it can accidentally cut a multibyte encoding unless additional validity conditions forbid that step.

Figure~\ref{fig:ch81-context} separates the preserved context from the replaced region.
Unlike a finite-state transducer, repeated rewriting feeds an output back as a new input and can iterate without bound.
An arbitrary string-rewrite relation should therefore not be assumed to admit one finite-state compilation merely because each individual left and right side is finite.

\begin{figure}[htbp]
\centering
\begin{tikzpicture}[node distance=15mm]
\node[basic block] (a) {prefix u | matched l | suffix v};
\node[basic block,below=of a] (b) {prefix u | replacement r | suffix v};
\draw[cfg edge] (a)--node[right,font=\small]{one rule occurrence}(b);
\end{tikzpicture}
\caption{A string rewrite preserves its surrounding context but can change all later offsets.}
\label{fig:ch81-context}
\end{figure}

\section{Work an overlap before claiming a canonical answer}
\label{sec:ch81-overlap}
Take the single rule \code{aa}$\to$\code{b} and input \code{aaa}.
Rewriting the left occurrence yields \code{ba}; rewriting the right yields \code{ab}.
Both are normal forms because neither contains \code{aa}.
The system terminates by length decrease but does not give a unique normal form.
Table~\ref{tab:ch81-overlap} computes the alternatives.

\begin{table}[htbp]
\caption{Length-decreasing rules can still produce incompatible normal forms.}
\label{tab:ch81-overlap}
\centering
\begin{tabular}{llll}
\toprule
Input & Occurrence start & Result & Further step \\
\midrule
aaa & 0 & ba & none \\
aaa & 1 & ab & none \\
\bottomrule
\end{tabular}
\end{table}

Choosing leftmost occurrence makes an implementation deterministic, but it does not change this fact about the underlying relation.
It defines one normalization policy, which may be acceptable for a text-processing tool if documented.
For canonicalization used as an equality decision, the distinction is essential: two convertible strings can otherwise receive unequal selected normal forms.
If a compiler uses normalized syntax as a cache key, unstable rule ordering can also alter cache behavior without changing the source program.

\section{Distinguish termination, normalization, and confluence}
\label{sec:ch81-properties}
A relation is \term{terminating} when it has no infinite forward reduction.
It is weakly normalizing on an element when at least one path reaches a normal form; other paths may diverge.
It is \term{confluent} when every pair of descendants of a common ancestor can be joined by further forward reductions.
Local confluence asks this joining property only for pairs of immediate successors.
These are properties of relations, not of the names of their algorithms~\cite{ohlebusch-2002-rewriting}.

Figure~\ref{fig:ch81-join} shows the confluence obligation.
The bottom node need not be a normal form.
Confluence implies that an element cannot have two different normal forms, but it does not promise that a normal form exists or that an arbitrary strategy reaches one.
Termination plus confluence gives existence and uniqueness for every element, which is the combination many canonicalizers want.

\begin{figure}[htbp]
\centering
\begin{tikzpicture}[node distance=14mm and 17mm]
\node[basic block] (a) {a};
\node[basic block,below left=of a] (b) {b};\node[basic block,below right=of a] (c) {c};
\node[basic block,below=of a,yshift=-24mm] (d) {join d};
\draw[cfg edge] (a)--(b);\draw[cfg edge] (a)--(c);
\draw[cfg edge,dashed] (b)--(d);\draw[cfg edge,dashed] (c)--(d);
\end{tikzpicture}
\caption{Confluence requires descendants of the same ancestor to admit a common descendant.}
\label{fig:ch81-join}
\end{figure}

Under termination, local confluence suffices for confluence; the termination premise cannot be dropped from that implication.
Chapter~\ref{ch:rewrite-confluence} develops critical overlaps and the exact premises needed to turn local checks into a global conclusion.
For a rewrite-based optimizer, semantic validity is another independent property: even a terminating confluent system can consistently compute the wrong program if one rule is unsound.

\section{Implement an explicit occurrence policy}
\label{sec:ch81-algorithm}
Algorithm~\ref{alg:ch81-string} selects the smallest start offset, breaking ties by declaration order.
Empty left sides are excluded so every match has a nonempty extent and occurrence enumeration remains simple.
An empty right side is allowed and deletes the matched substring.
The companion returns both the current string and a Boolean saying whether it actually reached a normal form.

\begin{algorithm}[htbp]
\caption{Normalize strings with a documented leftmost rule-order strategy.}
\label{alg:ch81-string}
\begin{pseudocode}
\Require finite string rules with nonempty left sides; input; nonnegative fuel
\Ensure selected normal form, or intermediate string with incomplete status
\While{fuel remains}
\State enumerate matching occurrences and their rule indices
\If{none exist} \State return current string and complete \EndIf
\State select minimum start offset, then minimum rule index
\State replace exactly that occurrence; decrement fuel
\EndWhile
\State test whether any occurrence remains; report completeness accordingly
\end{pseudocode}
\end{algorithm}

For $R$ rules, current length $N$, and maximum left-side length $P$, naive occurrence checking can cost $O(RNP)$ per step, plus replacement copying and occurrence-list storage.
For $S$ steps whose intermediate strings have length at most $N_{\max}$, this gives $O(SRN_{\max}P)$ matching work under direct comparison, not a bound on arbitrary rewrite termination.
An index can reduce matching cost, but length-changing replacement requires updating or rebuilding the affected index.
The small implementation favors an inspectable strategy over a production matcher.

On rules \code{ab}$\to$\code{a} and \code{aa}$\to$\code{a}, input \code{ababa} reduces to \code{aaba}, then \code{aba}, then \code{aa}, then \code{a}.
Every step shortens the string by one, so the four-step trace also matches the length potential from five to one.
By contrast, \code{a}$\to$\code{aa} grows forever from \code{a}; fuel exhaustion records an intermediate result, not a proof of termination or nontermination.

\section{Certify the trace, not an unrestricted equivalence claim}
\label{sec:ch81-certificate}
A rewrite certificate records the rule identity, occurrence, and intermediate or reconstructible state.
Algorithm~\ref{alg:ch81-replay} verifies each step against the actual current string.
Offsets must be interpreted after all previous replacements, not against the original input.
Replay establishes reachability under those rules; claiming semantic equivalence additionally requires that the rules are known sound for the object language.

\begin{algorithm}[htbp]
\caption{Replay a string-rewrite certificate against changing coordinates.}
\label{alg:ch81-replay}
\begin{pseudocode}
\Require original string, trusted rule table, ordered rule/offset events
\Ensure verified descendant or first invalid event
\ForAll{events $(r,i)$ in order}
\State check rule index and current substring extent
\State require the current substring at $i$ to equal rule $r$'s left side
\State replace that current occurrence with its right side
\EndFor
\State return reconstructed descendant
\end{pseudocode}
\end{algorithm}

Replay costs the sum of match-verification and replacement costs along the supplied trace.
It avoids searching for a redex but cannot avoid producing a large final output if the certificate expands the string.
The term engine in Chapter~\ref{ch:term-matching-unification} implements the analogous position-based replay on trees.
Production proof objects should also identify the rule-set version and semantic assumptions so an old trace cannot silently be checked against a different rule table.

\section{Use string rewriting only where token boundaries justify it}
\label{sec:ch81-compilers}
A textual rule replacing \code{x+0} with \code{x} ignores binding, precedence, comments, overloaded operators, and evaluation effects.
Even matching exact source bytes is not enough to justify a program transformation.
Token-string rewriting is useful for controlled macro expansion, instruction-sequence peepholes, or symbolic protocols when their state and boundary contracts are explicit.
Tree and graph representations expose more semantic structure for general compiler optimization.

Table~\ref{tab:ch81-levels} contrasts the carriers.
A graph redex can have many incoming uses and does not correspond to one contiguous textual interval.
A term position selects one occurrence in an unfolded tree, while a shared graph node can represent several such occurrences.
The later graph chapters build that distinction into update algorithms instead of treating pointers as a mere implementation detail.

\begin{table}[htbp]
\caption{Different rewrite carriers require different notions of occurrence and context.}
\label{tab:ch81-levels}
\centering
\begin{tabular}{lll}
\toprule
Carrier & Occurrence & Hidden obligation \\
\midrule
byte string & byte interval & encoding and source syntax \\
token string & token interval & binding and evaluation order \\
term tree & child-index path & substitution and types \\
shared term graph & node and graph match & aliasing and sharing semantics \\
compiler CFG & blocks/edges/value uses & dominance, effects, control validity \\
\bottomrule
\end{tabular}
\end{table}

\section{Reason about the induced equality rather than one chosen trace}
\label{sec:ch81-conversion}
The directed relation and its symmetric conversion closure serve different purposes.
A directed normalizer chooses a representation; the conversion relation describes the equations generated by the rules.
For aa$\to$b, both ba and ab are convertible through their common predecessor aaa even though neither reduces to the other.
Consequently a test that compares their selected normal forms would incorrectly declare them inequivalent if it assumed confluence without establishing it.
This is not a defect in the leftmost implementation: it is an unsupported use of that implementation as a decision procedure.

The distinction appears in compiler caches.
Suppose two optimization paths produce semantically equal expressions, but the cache key is the printed output of a nonconfluent simplifier.
The keys may differ and reduce cache reuse without causing a miscompilation.
If the compiler instead treats unequal keys as proof that the expressions are semantically different, it has crossed from a performance issue into an invalid reasoning rule.
Sound positive evidence and complete equality decisions are different interface promises.

A deterministic terminating strategy can still be useful in a nonconfluent system.
Its output is reproducible when rule order and occurrence policy are fixed.
To obtain idempotence, the strategy must actually stop only when none of its own admitted steps remains; an arbitrary one-pass scan can expose a redex behind the scan point and fail that property.
Even idempotence is weaker than canonical equality: repeatedly returning the same chosen normal form for each input does not guarantee that convertible inputs share a result.

\section{Model conditional and contextual string rules explicitly}
\label{sec:ch81-conditional}
Real textual transformations often require conditions such as being outside a string literal, having a token boundary on each side, or not crossing a comment.
Those conditions belong to the relation.
For instance, removing repeated spaces may preserve a token stream between identifiers while changing the contents of a quoted literal.
The same byte pattern can therefore be a redex in one context and ordinary data in another.
A rule engine that scans raw bytes must either carry lexical state or operate on a representation where the relevant contexts are already distinguished.

Boundary predicates can change after replacement.
If deleting a separator joins two identifiers, a later match based on the old tokenization is stale.
A safe engine rescans the affected lexical neighborhood or uses a representation with explicit token objects and a verified retokenization rule.
Recording old byte offsets without the source version is insufficient because all later offsets can shift after a length-changing step.
The certificate checker must apply each event against the current state, as Algorithm~\ref{alg:ch81-replay} requires.

Termination analysis must include the conditions too.
An unconditional length decrease is sufficient even when conditions restrict which steps occur, because restricting a terminating relation cannot create an infinite path.
But adding a new conditional expansion rule is not harmless merely because it rarely triggers.
The proof must establish a well-founded decrease for every permitted occurrence, or use a phase-specific argument that prevents the expansion from being repeated indefinitely.
Counting observed successful runs is useful testing evidence, not a universal termination proof.

\section{Derive uniqueness from the joining property}
\label{sec:ch81-uniqueness}
Suppose a has two normal-form descendants n and m and the relation is confluent.
Confluence supplies a common descendant d with $n\to^*d$ and $m\to^*d$.
Since n and m are normal, neither can make a positive step.
Both paths therefore contain zero steps, so n=d=m.
This proves uniqueness whenever a normal form exists without assuming termination of every reduction path.

Existence needs another argument.
For a terminating finitely presented rewrite system with a complete effective redex search, repeatedly taking an available step cannot continue forever and eventually stops at a normal form.
Combining that existence argument with confluence yields a unique result for every starting object.
The implementation still needs a complete search and correct matching; a missed redex can make it report a premature normal form even when the mathematics is sound.

These proof boundaries suggest separate regression tests.
Replay checks that each reported step is permitted.
A final full redex scan checks normality.
Comparing different strategies on finite generated inputs can expose nonconfluence or implementation defects, but agreement on the tested inputs does not establish global confluence.
Finally, an independent source evaluator checks that the selected rewriting relation preserves the intended program behavior.
No one of these tests subsumes the others.

\section{Exercises}
\label{sec:ch81-exercises}
\begin{enumerate}
\item\label{ex:ch81-overlap} Enumerate all normal forms reachable from \code{aaa} using \code{aa}$\to$\code{b}.
\item\label{ex:ch81-trace} Replay the four-step \code{ababa} example, recording the offset and rule chosen at each step.
\item\label{ex:ch81-fuel} Apply \code{a}$\to$\code{aa} for three steps from \code{a}; state the result and what the fuel status does and does not establish.
\item\label{ex:ch81-properties} Explain why a length-decreasing rule collection proves termination but does not, by itself, prove confluence or semantic correctness.
\end{enumerate}

\section{Further reading}
\label{sec:ch81-reading}
Ohlebusch develops abstract reduction, termination, and confluence with their necessary distinctions~\cite{ohlebusch-2002-rewriting}.
Chapter~\ref{ch:finite-state-transducers} provides the complementary finite-state relation model; the two models should not be conflated.
